\documentclass[a4paper,12pt]{report}

% === Encoding and Language ===
\usepackage[utf8]{inputenc}
\usepackage[T1]{fontenc}
\usepackage[english]{babel}

% === Mathematics and Figures ===
\usepackage{amsmath, amssymb, amsthm}
\usepackage{graphicx}
\graphicspath{{../results/}}
\usepackage{float}
\usepackage{booktabs}
\usepackage{multirow}
\usepackage{array}
\usepackage{algorithm}
\usepackage{algpseudocode}

% === Page Layout and References ===
\usepackage[margin=1in]{geometry}
\usepackage{parskip}
\usepackage{hyperref}
\hypersetup{
    colorlinks=true,
    linkcolor=blue,
    citecolor=red,
    urlcolor=magenta,
}

% === Bibliography ===
\usepackage{natbib}
\bibliographystyle{plainnat}

% === Custom Commands ===
\newcommand{\thesistitle}{%
  Hierarchical Minimum Variance Allocation: Nonlinear Shrinkage,
  Vol-Balanced Trees, Black-Litterman Returns, and Kalman-Filter
  Smoothing for Large-Cap Equity Portfolios}
\newcommand{\authorname}{Kiril Ivanov}
\newcommand{\supervisor}{Sanders Barendse}
\newcommand{\submissiondate}{\today}

\newtheorem{proposition}{Proposition}
\newtheorem{remark}{Remark}

\begin{document}

% === Title Page ===
\begin{titlepage}
  \centering
  {\scshape\LARGE University of Amsterdam \par}
  \vspace{1cm}
  {\scshape\Large BSc Econometrics and Data Science Thesis\par}
  \vspace{2cm}
  {\huge\bfseries \thesistitle \par}
  \vspace{2cm}
  {\Large \authorname \par}
  \vspace{1cm}
  {\large Supervisor: \supervisor \par}
  \vfill
  {\large \submissiondate \par}
\end{titlepage}

% === Abstract ===
\newpage
\chapter*{Abstract}
\addcontentsline{toc}{chapter}{Abstract}

This thesis proposes and evaluates the \textbf{Hierarchical Minimum Variance
Allocator} (HMVA), a six-stage modification of Hierarchical Risk Parity (HRP)
that replaces sample covariance with analytical nonlinear shrinkage (NLS) applied
to exponentially weighted returns (Stages 1--2), adds Black-Litterman expected
returns based on a skip-one-month momentum signal (Stage 3), replaces the
agglomerative correlation tree with a vol-balanced top-down splitting criterion
(Stage 4), uses Sharpe-ratio-weighted bisection instead of variance-based
bisection (Stage 5), and applies a Kalman-filter weight smoother to control
turnover and noise-chasing in weight updates (Stage 6).

An ablation study on the top 100 S\&P~500 stocks by market capitalisation, using
a fully point-in-time CRSP universe from 2002 to 2024 and a two-year estimation
window ($T = 504$ trading days), documents the Sharpe contribution of each stage.
HMVA achieves an annualised Sharpe ratio of 0.809, compared to 0.574 for standard
HRP (sample covariance), and 0.477 for the $1/N$ equal-weight benchmark. The
cumulative improvement over HRP is $+0.235$ Sharpe units, decomposed as:
$+0.019$ from NLS and EWMA, $+0.042$ from the vol-balanced tree, $+0.068$ from
Black-Litterman returns with Sharpe bisection, and $+0.105$ from the Kalman-filter
smoother. HMVA simultaneously achieves a lower maximum drawdown ($-38.1\%$) and
a higher Sortino ratio (1.076) than HRP ($-41.8\%$, 0.727) and equal-weight
($-55.2\%$, 0.596). A cost-robustness sweep across 20 cells (four lookback windows
$\times$ five transaction-cost levels from 0 to 20~bps) confirms that HMVA
outperforms equal-weight in every cell, with a mean Sharpe advantage of 0.301.
Crisis-period analysis shows HMVA substantially outperforms during the Global
Financial Crisis (annualised Sharpe $-0.31$ versus $-0.88$ for HRP and $-0.89$
for equal-weight) and the COVID-19 crash; performance is mixed during the
dot-com trough and the 2022 rate-hike cycle.

% === Table of Contents ===
\tableofcontents
\listoffigures
\listoftables

% =============================================================================
\chapter{Introduction}
\label{ch:intro}

The construction of well-diversified equity portfolios remains a central problem
in quantitative asset management. Since \citet{markowitz1952}, the dominant
framework has been mean-variance optimisation: estimate expected returns $\mu$ and
a covariance matrix $\Sigma$, then allocate capital to the portfolio on the
efficient frontier. In practice both inputs must be estimated from finite samples,
and the error in $\hat{\Sigma}$ enters the portfolio through matrix inversion in
the formula for the global minimum-variance portfolio:
\begin{equation}
  w^{\text{GMV}} = \frac{\hat{\Sigma}^{-1}\mathbf{1}}{\mathbf{1}^{\top}\hat{\Sigma}^{-1}\mathbf{1}}.
  \label{eq:intro-gmv}
\end{equation}
Because inversion amplifies small eigenvalues, near-zero eigenvalues of the sample
covariance --- which arise whenever $N$ is not much smaller than $T$ --- generate
extreme and unstable portfolio weights \citep{michaud1989}.

Two families of solutions have been proposed. The first regularises the
covariance estimate itself. \citet{ledoitwolf2004} derived the analytically
optimal linear shrinkage intensity toward a structured target, and
\citet{ledoitwolf2020} extended this to eigenvalue-by-eigenvalue nonlinear
shrinkage (NLS). The second family avoids matrix inversion entirely.
\citet{lopezdeprado2016} proposed Hierarchical Risk Parity (HRP): build a
hierarchy from the asset correlation matrix and allocate risk recursively through
the tree using variance-based bisection, never inverting $\Sigma$.

HRP's non-amplifying property raises a natural follow-on question: if estimation
error is not blown up by inversion, does the quality of the covariance estimate
still matter? Standard HRP uses the sample covariance in both the correlation
distance matrix (which determines tree structure) and the bisection weights
(which determine how capital is split at each node). A better estimate of $\Sigma$
should improve both channels simultaneously. At the same time, HRP's purely
risk-based allocation ignores expected returns entirely, which \citet{demiguel2009}
showed to be a high hurdle only occasionally cleared by sophisticated optimisers.
Incorporating a regularised return signal through Black-Litterman estimation
\citep{blacklitterman1992} offers a principled way to tilt the allocation without
the instability of raw sample means.

Beyond the covariance and return inputs, the tree-construction criterion and the
weight-update dynamics of HRP also offer room for improvement. Standard HRP builds
the tree by agglomerative clustering, which is designed to maximise correlation
separation but ignores the risk balance across branches. A top-down tree that
directly minimises a blended vol-balance and correlation objective produces a more
risk-equitable hierarchy. Once target weights are computed, applying them naively
every period creates excessive turnover driven by signal noise; a Kalman-filter
smoother that conditions weight updates on the signal-to-noise ratio of the return
estimate can substantially improve the risk-adjusted performance of the
continuously-rebalanced portfolio.

The present thesis pursues a single contribution: the Hierarchical Minimum Variance
Allocator (HMVA), a six-stage pipeline that systematically addresses these known
failure modes of standard HRP. A fully point-in-time backtest on the top-100
S\&P~500 stocks from 2002 to 2024 --- with no look-ahead bias and correct
survivorship-bias avoidance --- provides the empirical evaluation. A stage-by-stage
ablation study isolates the Sharpe contribution of each component, avoiding the
common pitfall of attributing all performance to the full model without
understanding which parts drive the result.

The remainder of this thesis is organised as follows. Chapter~\ref{ch:litreview}
reviews the relevant portfolio theory, covariance estimation, and HRP literature.
Chapter~\ref{ch:methodology} derives all estimators and the full HMVA pipeline.
Chapter~\ref{ch:data} describes the CRSP data and the point-in-time universe
construction. Chapter~\ref{ch:results} presents the ablation study and main
empirical results. Chapter~\ref{ch:robustness} reports the cost and parameter
robustness analyses. Chapter~\ref{ch:crisis} presents the crisis-period analysis.
Chapter~\ref{ch:discussion} discusses the findings, limitations, and directions
for future work. Chapter~\ref{ch:conclusion} concludes.

% =============================================================================
\chapter{Literature Review}
\label{ch:litreview}

\section{Mean-variance optimisation and its estimation problem}
\label{sec:lit-mv}

\citet{markowitz1952} established that an investor who cares only about the first
two moments of portfolio returns should hold a portfolio on the efficient frontier.
The global minimum-variance (GMV) portfolio, which minimises variance subject to
full investment and long-only constraints, requires no estimate of expected returns
and is defined by~\eqref{eq:intro-gmv}.

The sample covariance $S = (T-1)^{-1}X^{\top}X$ (where $X$ is the $T\times N$
demeaned return matrix) is the natural estimator of $\Sigma$, but when $N$ is not
small relative to $T$ it is unreliable. By the Marchenko-Pastur law
\citep{marchenkopastur1967}, the sample eigenvalues $\{d_i\}$ are systematically
biased: large eigenvalues are inflated and small ones are deflated. Inverting $S$
amplifies the small (deflated) eigenvalues, causing the GMV portfolio to load
heavily onto directions dominated by noise. \citet{michaud1989} termed this
``optimisation error maximisation.'' \citet{chopraziemba1993} showed that errors
in the covariance matrix are approximately ten times less damaging than errors in
expected returns, yet still materially degrade portfolio performance.

\citet{demiguel2009} provided the starkest empirical benchmark: across 14
strategies and seven datasets, no sophisticated optimiser consistently outperformed
the naive $1/N$ equal-weight portfolio once estimation error was accounted for.
This finding defines the minimum hurdle for any new strategy.

\section{Covariance regularisation}
\label{sec:lit-cov}

\subsection{Ledoit-Wolf linear shrinkage}

\citet{ledoitwolf2004} proposed shrinking the sample covariance linearly toward
a structured target $F$:
\begin{equation}
  \hat{\Sigma}_{\text{LW}} = \alpha^{\star} F + (1-\alpha^{\star}) S,
  \label{eq:lw}
\end{equation}
where $F$ is the constant-correlation covariance (all pairwise correlations equal
their cross-sectional average) and $\alpha^{\star}$ minimises the expected
Frobenius loss $\mathbb{E}\|\hat{\Sigma}-\Sigma\|_F^2$. The intensity
$\alpha^{\star}$ is computed in closed form and increases automatically when
estimation is unreliable. The estimator is positive-definite for all $N, T \geq 2$.

\citet{ledoitwolf2020} showed that linear shrinkage toward a single target is
suboptimal because each eigenvalue has a different optimal shrinkage amount.
Their analytical nonlinear shrinkage (NLS) estimator replaces each sample
eigenvalue $d_i$ with an oracle-optimal correction derived from the Stieltjes
transform of the sample spectral distribution:
\begin{equation}
  d_i^{\star} = \frac{d_i}{\bigl|1 - (N/T) - (N/T)\, d_i\, \tilde{m}_F(d_i)\bigr|^2},
  \label{eq:nls}
\end{equation}
while the eigenvectors are unchanged. NLS achieves strictly smaller expected
Frobenius loss than any linear shrinkage estimator as $N, T \to \infty$ with
$N/T \to c \in (0, \infty)$.

\subsection{EWMA covariance}

\citet{riskmetrics1996} introduced the exponentially weighted moving average
(EWMA) covariance, which assigns geometrically declining weights to past
observations:
\begin{equation}
  \hat{\Sigma}^{\text{EWMA}} = \sum_{t=0}^{T-1} w_t \, r_t r_t^{\top},
  \qquad
  w_t = \frac{(1-\lambda)\lambda^{T-1-t}}{\sum_{s=0}^{T-1}(1-\lambda)\lambda^s},
  \label{eq:ewma}
\end{equation}
with decay parameter $\lambda \in (0,1)$. EWMA makes the covariance estimate
respond faster to volatility clustering than an equally weighted rolling window,
which is particularly valuable around market stress events.

\section{POET: factor-based thresholding}
\label{sec:lit-poet}

\citet{fanliamincheva2013} proposed POET (Principal Orthogonal complEment
Thresholding), which exploits the approximate factor model $r_t = Bf_t + u_t$.
The covariance decomposes as $\Sigma = B\Sigma_f B^{\top} + \Sigma_u$ and POET
estimates it in two steps: extract the $K$ leading eigen-pairs of $S$ as the
factor component, then adaptively threshold the residual $\hat{\Sigma}_u$ to
enforce sparsity. POET is consistent under a factor-plus-sparse true structure;
when no such structure exists, the forced decomposition introduces additional
estimation error.

\section{Hierarchical Risk Parity}
\label{sec:lit-hrp}

\citet{lopezdeprado2016} proposed HRP to avoid the matrix-inversion bottleneck
of mean-variance. The algorithm proceeds in three stages:
\begin{enumerate}
  \item Compute a correlation-distance matrix $D_{ij} = \sqrt{(1-\hat{\rho}_{ij})/2}$
        and build a hierarchical clustering tree by agglomerative linkage.
  \item Reorder assets so that correlated assets are adjacent in the leaf sequence
        (quasi-diagonalisation).
  \item Recursively bisect the tree: at each internal node, split the weight between
        left and right subtrees inversely proportional to their equal-weight cluster
        variances (variance bisection).
\end{enumerate}

By not inverting $\Sigma$, HRP is immune to the Marchenko-Pastur amplification
effect. However, the covariance estimate still enters through both the distance
matrix (governing tree structure) and the bisection weights (governing capital
allocation), so a better estimate of $\Sigma$ should improve both. Separately,
the agglomerative tree optimises for correlation separation, not risk balance across
branches, which can leave substantial variance imbalance between subtrees.
\citet{molyboga2020} proposed modifications to the bisection step to incorporate
return estimates; HMVA builds on this idea while also replacing the tree-construction
criterion.

\section{Black-Litterman return estimation}
\label{sec:lit-bl}

\citet{blacklitterman1992} introduced a Bayesian framework for blending equilibrium
expected returns (the CAPM-implied prior $\pi = \delta\Sigma w_{\text{mkt}}$) with
investor views $Q$ under uncertainty $\Omega$, yielding a posterior return vector:
\begin{equation}
  \mu_{\text{BL}} = \bigl[(\tau\hat{\Sigma})^{-1} + P^{\top}\Omega^{-1}P\bigr]^{-1}
                   \bigl[(\tau\hat{\Sigma})^{-1}\pi + P^{\top}\Omega^{-1}Q\bigr],
  \label{eq:bl}
\end{equation}
with $P$ the view matrix. The BL posterior is more stable than the sample mean
and less sensitive to noise, because the prior regularises extreme sample means
toward the equilibrium.

\section{Research gap}
\label{sec:lit-gap}

Existing work compares HRP against mean-variance \citep{lopezdeprado2016},
modifies the bisection step \citep{molyboga2020}, or evaluates covariance
estimators in mean-variance portfolios \citep{ledoitwolf2004, fanliamincheva2013}.
No published study has (i)~combined NLS covariance, EWMA pre-weighting, BL
expected returns, a vol-balanced top-down tree, Sharpe-ratio bisection, and a
Kalman-filter smoother into a single HRP pipeline; (ii)~evaluated this pipeline
on a fully point-in-time S\&P~500 universe with no survivorship bias; or
(iii)~provided a stage-by-stage ablation that isolates the marginal Sharpe
contribution of each component. The present thesis fills this gap.

% =============================================================================
\chapter{Methodology}
\label{ch:methodology}

\section{Baseline portfolio strategies}
\label{sec:meth-baselines}

\subsection{Standard Hierarchical Risk Parity (HRP)}

At each rebalance date, compute the sample covariance $S$ over the lookback window
$[t-T, t-1]$, derive the correlation-distance matrix
$D_{ij} = \sqrt{(1-\hat{\rho}_{ij})/2}$, and build a clustering tree using
single-linkage agglomeration. Weights are assigned by recursive variance bisection:
\begin{equation}
  \alpha_L = 1 - \frac{\tilde{\sigma}_L^2}{\tilde{\sigma}_L^2 + \tilde{\sigma}_R^2},
  \qquad \alpha_R = 1 - \alpha_L,
  \label{eq:hrp-bisect}
\end{equation}
where $\tilde{\sigma}_C^2 = \mathbf{e}^{\top}S_C\mathbf{e}/n_C^2$ is the
equal-weight variance of cluster $C$. Standard HRP uses the sample covariance
throughout. This is the primary baseline in the ablation study.

\subsection{Equal-weight (EW)}

The $1/N$ equal-weight portfolio assigns weight $1/N$ to each asset at each
rebalance without any estimation. \citet{demiguel2009} demonstrated that EW is
difficult to beat on a risk-adjusted basis, making it the primary statistical
baseline.

\subsection{Market-cap weighted (SPY-K)}

Weights are proportional to market capitalisation among the top-100 stocks at
each rebalance date. This benchmark approximates the S\&P~100 index allocation.

\subsection{Mean-variance utility (MVO)}

Maximises the long-only utility $w^{\top}\mu - (\gamma/2)w^{\top}\Sigma w$ with
risk aversion $\gamma = 2.5$ and Black-Litterman posterior expected returns via
SLSQP. This serves as a reference for fully parametric optimisation.

\section{Covariance estimators}
\label{sec:meth-cov}

\subsection{Sample covariance}

$S = (T-1)^{-1}X^{\top}X$ where $X$ is the demeaned return matrix. Used as input
to the baseline HRP strategy.

\subsection{EWMA pseudo-returns preprocessing}

Given a $T \times N$ return matrix $R$ and EWMA half-life $h$
($\lambda = 0.5^{1/h}$), the pseudo-return matrix is:
\begin{equation}
  \tilde{R}_{s,\cdot} = \sqrt{w_s \cdot T}\; R_{s,\cdot},
  \qquad w_s = \frac{(1-\lambda)\lambda^{T-1-s}}{\sum_{j=0}^{T-1}(1-\lambda)\lambda^j},
  \label{eq:ewma-pseudo}
\end{equation}
so that $T^{-1}\tilde{R}^{\top}\tilde{R} = \hat{\Sigma}^{\text{EWMA}}$. Passing
$\tilde{R}$ to NLS allows the estimator to inherit exponential weighting without
modifying its internals. With $h = 21$ trading days ($\lambda \approx 0.967$),
the effective weight on observations decays to half within one month.

\subsection{Analytical nonlinear shrinkage (NLS)}

Applied to $\tilde{R}$ from~\eqref{eq:ewma-pseudo}, equation~\eqref{eq:nls}
yields $\hat{\Sigma}^{\text{NLS}}$. This matrix is positive-definite, well-conditioned,
and corrects the systematic eigenvalue bias of the sample covariance at any $N/T$
ratio. Its eigenvectors --- the principal directions of common variation among assets
--- are identical to those of the sample covariance, so the tree structure inherits
the same information content while using more accurate eigenvalue magnitudes in the
bisection weights.

\section{The HMVA: Hierarchical Minimum Variance Allocator}
\label{sec:meth-hmva}

Table~\ref{tab:hmva-stages} summarises the six stages. Each directly addresses
a known failure mode of standard HRP, as established in the literature or as
identified empirically in the ablation study.

\begin{table}[H]
  \centering
  \caption{HMVA pipeline: stages, operations, and failure modes addressed.}
  \label{tab:hmva-stages}
  \begin{tabular}{@{}clp{6cm}r@{}}
    \toprule
    Stage & Operation & Failure mode addressed & $\Delta$Sharpe \\
    \midrule
    1 & EWMA pseudo-returns ($h=21$) & Equal weighting ignores vol.\ clustering & \multirow{2}{*}{$+0.019$} \\
    2 & NLS covariance               & Sample $\hat{\Sigma}$ has biased eigenvalues & \\
    3 & Black-Litterman $\hat{\mu}$  & Pure-risk allocation ignores returns & \multirow{2}{*}{$+0.068$} \\
    5 & Sharpe-ratio bisection        & Variance bisection ignores risk-adj.\ return & \\
    4 & Vol-balanced top-down tree   & Corr.\ tree ignores branch risk balance & $+0.042$ \\
    6 & Kalman-filter weight smoother & Naive updates chase noise, inflate turnover & $+0.105$ \\
    \bottomrule
    \multicolumn{3}{l}{\small Total improvement over HRP baseline:} & $\mathbf{+0.235}$ \\
  \end{tabular}
\end{table}

\subsection{Stage 1: EWMA pseudo-returns}

Compute EWMA weights $\{w_s\}_{s=0}^{T-1}$ with half-life $h = 21$ and form
the pseudo-return matrix $\tilde{R}$ via~\eqref{eq:ewma-pseudo}. This ensures
that the NLS estimator in Stage~2 operates on an exponentially weighted
covariance without modification to the estimator's internals.

\subsection{Stage 2: NLS covariance}

Apply the Ledoit-Wolf (2020) analytical nonlinear shrinkage formula~\eqref{eq:nls}
to $\tilde{R}$, yielding $\hat{\Sigma}^{\text{NLS}}$. The resulting matrix is
the primary covariance input to all subsequent stages.

\subsection{Stage 3: Black-Litterman expected returns}

Using the \emph{raw} (not EWMA-weighted) return window to avoid double-discounting
of the mean signal:
\begin{enumerate}
  \item Compute the skip-one-month momentum signal by averaging returns over
        $[0, T - 21)$: $Q = (T-21)^{-1}\sum_{t < T-21} r_t$.
        The most recent 21 trading days are skipped to avoid the reversal effect
        in short-horizon momentum.
  \item Equilibrium prior: $\pi = \delta\,\hat{\Sigma}^{\text{NLS}}\,\mathbf{1}/N$,
        where $\mathbf{1}/N$ is the equal-weight proxy for the market portfolio
        and $\delta = 2.5$ is the risk-aversion coefficient.
  \item View matrix $P = I_N$ (one absolute view per asset), view vector $Q$,
        view uncertainty $\Omega = \tau\hat{\Sigma}^{\text{NLS}}$ with $\tau = 0.05$.
        With $\tau = 0.05$ the prior receives $\approx 5\%$ weight in the posterior.
  \item BL posterior via~\eqref{eq:bl}:
        $\mu_{\text{BL}} \leftarrow [\hat{\Sigma}^{-1}/\tau + \Omega^{-1}]^{-1}
        [\hat{\Sigma}^{-1}\pi/\tau + \Omega^{-1}Q]$, which simplifies to
        $\pi + \hat{\Sigma}(\hat{\Sigma}+\Omega)^{-1}(Q-\pi)$ when $P = I$.
        Extreme sample means are pulled toward the CAPM equilibrium, reducing
        sensitivity to short-horizon estimation noise.
\end{enumerate}

\remark{BL parameters $\delta$ and $\tau$ cancel analytically in the posterior
formula when $P = I$ and $\Omega = \tau\hat{\Sigma}$. Specifically, the posterior
depends only on $\hat{\Sigma}(\hat{\Sigma}+\Omega)^{-1} = I/(1+\tau)$, which is
a constant scaling. Sensitivity sweeps confirm that varying these parameters has
no empirical effect on out-of-sample Sharpe, consistent with the analytical result.}

\subsection{Stage 4: Vol-balanced top-down tree}

Standard HRP builds the tree by agglomerative clustering of correlation distances.
This criterion groups highly correlated assets together but pays no attention to
whether the resulting branches are risk-balanced: two branches may differ
substantially in total volatility, causing the bisection weights to concentrate
capital in the low-volatility branch.

HMVA replaces agglomerative clustering with a \emph{top-down greedy binary split}
that minimises a blended objective at each node. For a bipartition $(A, B)$ of
a node's assets, the split cost is:
\begin{equation}
  f(A, B;\,\lambda) = \lambda\,\mathrm{VB}(A,B) + (1-\lambda)\,\rho(A,B),
  \label{eq:split-obj}
\end{equation}
where
\begin{equation}
  \mathrm{VB}(A,B) = \frac{|\mathrm{vol}(A) - \mathrm{vol}(B)|}{\mathrm{vol}(A) + \mathrm{vol}(B)}
\end{equation}
is the equal-weight volatility imbalance (0 = perfectly balanced, 1 = all risk in
one branch) and
\begin{equation}
  \rho(A,B) = \frac{\sum_{i \in A, j \in B}\hat{\Sigma}_{ij}}
                   {\sqrt{\bigl(\sum_{i,j \in A}\hat{\Sigma}_{ij}\bigr)
                          \bigl(\sum_{i,j \in B}\hat{\Sigma}_{ij}\bigr)}}
\end{equation}
is the equal-weight inter-cluster correlation (computed directly from the
covariance submatrices without inversion). The blend parameter $\lambda = 0.25$
assigns 25\% weight to vol-balance and 75\% to correlation minimisation,
producing a tree that is more risk-equitable than pure correlation clustering
while retaining most of its diversification structure.

\textbf{Implementation.} For clusters with $\leq 10$ assets, an exhaustive search
over all $2^{n-1}/2$ bipartitions evaluates the exact minimum of~\eqref{eq:split-obj}.
For clusters with $> 10$ assets, a heuristic sorts assets by within-cluster
row-sum (a proxy for systematic risk) and evaluates all $n-1$ contiguous cuts
using 2-D prefix sums in $O(n^2)$ time.

\subsection{Stage 5: Sharpe-ratio bisection}

At each internal node with children $L$ and $R$, compute the cluster Sharpe proxy:
\begin{equation}
  \hat{S}(C) = \max\!\left(\frac{\bar{\mu}_{\text{BL},C}}{\sigma_{\text{EW}}(C)},\; 0\right),
  \label{eq:sharpe-proxy}
\end{equation}
where $\bar{\mu}_{\text{BL},C} = |C|^{-1}\sum_{i \in C}\mu_{\text{BL},i}$ is
the mean BL return and $\sigma_{\text{EW}}(C) = \sqrt{\mathbf{e}^{\top}\hat{\Sigma}_C\mathbf{e}/n_C^2}$
is the equal-weight cluster volatility. The weight allocation is:
\begin{equation}
  \alpha_L = \frac{\hat{S}(L)}{\hat{S}(L)+\hat{S}(R)}, \qquad \alpha_R = 1 - \alpha_L,
  \label{eq:sharpe-alloc}
\end{equation}
so that the branch with the higher risk-adjusted return receives more capital. If
both Sharpe proxies are zero or negative (which can occur in prolonged downtrends),
the split reverts to inverse-variance bisection~\eqref{eq:hrp-bisect}, ensuring
the allocator never applies negative weight to an asset.

The tree structure inherited from Stage~4 bounds the maximum allocation to any
single cluster, preventing the extreme concentration that unconstrained
mean-variance optimisation can produce while still incorporating the BL return
signal.

\subsection{Stage 6: Kalman-filter weight smoother}

Standard HRP applies the new target weights $w_{\text{raw}}$ mechanically at each
rebalance. When the BL signal is noisy --- for example when the estimation period
straddles a regime change --- this generates large weight swings that increase
realised transaction costs without improving expected performance. The Kalman-filter
smoother interpolates between the previous period's weights and the new target
according to a time-varying gain that reflects the current signal-to-noise ratio.

\textbf{Measurement noise} $R_t$ is proxied by the spectral entropy of
$\hat{\Sigma}^{\text{NLS}}$:
\begin{equation}
  R_t = \frac{-\sum_{i=1}^{N} p_i \log p_i}{\log N}, \qquad
  p_i = \frac{\lambda_i}{\sum_j \lambda_j},
  \label{eq:kf-r}
\end{equation}
where $\{\lambda_i\}$ are the eigenvalues of $\hat{\Sigma}^{\text{NLS}}$.
$R_t \in [0,1]$: it equals 1 when all eigenvalues are equal (maximally
uncertain covariance structure) and approaches 0 when a single eigenvalue
dominates. A high $R_t$ signals an ill-conditioned estimation environment; the
smoother should trust the new target less.

\textbf{Process noise} $Q_t$ is proxied by the relative change in the BL return
vector:
\begin{equation}
  Q_t = \frac{\|\mu_{\text{BL},t} - \mu_{\text{BL},t-1}\|}{\|\mu_{\text{BL},t-1}\| + \epsilon},
  \label{eq:kf-q}
\end{equation}
where $\epsilon$ is a small positive constant. A high $Q_t$ indicates that the
return signal has moved substantially, which warrants a larger weight update.

\textbf{Gain and smoothed weights:}
\begin{equation}
  K_t = \frac{Q_t}{Q_t + R_t}, \qquad
  w_t = (1-K_t)\,w_{t-1} + K_t\,w_{\text{raw},t}.
  \label{eq:kf-update}
\end{equation}
When the covariance structure is stable ($R_t$ high) and the return signal is
small ($Q_t$ low), $K_t$ is small and the portfolio barely moves, conserving
transaction capacity. When the signal changes substantially relative to structural
uncertainty, $K_t$ is large and the portfolio updates aggressively. The smoothed
weights are renormalised to sum to one after applying~\eqref{eq:kf-update}.

\section{Statistical inference}
\label{sec:meth-inference}

\subsection{Block-bootstrap Sharpe ratio test}

Following \citet{ledoitwolf2008}, a circular block bootstrap (block length $b=21$,
corresponding to one calendar month) is used to estimate the distribution of the
Sharpe ratio difference $\text{SR}^{\text{HMVA}} - \text{SR}^{B}$ under the null
of equality. The empirical one-sided $p$-value is the fraction of $B=2000$
bootstrap samples for which the bootstrapped difference exceeds the observed
difference. Block length 21 preserves the approximate monthly autocorrelation
structure of equity returns.

\subsection{Multiple testing correction}

All pairwise comparisons against the EW baseline are corrected for multiple testing.
The Holm-Bonferroni step-down procedure controls the family-wise error rate (FWER)
at the 5\% level. The Benjamini-Hochberg (BH) procedure controls the false discovery
rate (FDR) at 5\% and is more powerful when several strategies genuinely outperform.

% =============================================================================
\chapter{Data and Universe Construction}
\label{ch:data}

\section{CRSP data}
\label{sec:data-crsp}

The empirical experiment uses daily CRSP data in CIZ format with four primary fields:
\texttt{PERMNO} (unique stock identifier), \texttt{DlyCalDt} (calendar date),
\texttt{DlyClose} (split-adjusted closing price), and \texttt{DlyRet} (daily total
return). When \texttt{DlyRet} is available for a given stock-date pair it is used
directly; otherwise the return is computed from consecutive adjusted closing prices.
Daily log-returns are $r_{i,t} = \log(1 + \texttt{DlyRet}_{i,t})$.

\textbf{Delisting treatment.} Stocks that exit the panel mid-sample are handled as
follows: if CRSP already records a large negative return on the last observation
(likely capturing the actual delisting loss), the observation is kept as-is. If
the stock's last recorded return is in the ``normal trading'' range (between the
delisting return threshold and $+10\%$), a $-30\%$ return is imputed on the
first missing day, following \citet{shumway1997} who estimates the average
delisting loss at approximately 30\% for involuntary delistings.

\section{Point-in-time investment universe}
\label{sec:data-pit}

S\&P~500 membership is reconstructed from a historical constituents file that
records membership spells in range format: each row identifies a PERMNO and the
interval during which it was an S\&P~500 constituent. This prevents survivorship
bias: only stocks that were members of the S\&P~500 on the specific rebalance
date are eligible for inclusion in the portfolio.

At each 21-day rebalance date $t$, the investable universe $\mathcal{U}_t$ is
constructed as follows:
\begin{enumerate}
  \item Find the set of PERMNOs with active S\&P~500 membership on $t$ from the
        constituents file.
  \item Restrict to the top 100 by market capitalisation as of $t$, using the
        most recent available daily cap figure.
  \item Retain only those PERMNOs with complete non-\texttt{NaN} return data over
        the full estimation window $[t - T_{\max}, t - 1]$ where $T_{\max} = 504$
        days. Stocks that fail this criterion (typically recent IPOs or stocks
        with gaps due to trading halts) are excluded from all strategies at that
        rebalance.
  \item Estimate $\hat{\Sigma}$ on the lookback window $[t - T, t - 1]$ using the
        appropriate estimator for each strategy.
\end{enumerate}

The resulting universe $N_t \leq 100$ changes at every rebalance date, reflecting
actual index additions, deletions, and corporate events in real time. The backtest
period spans January 2002 to December 2024, with the first usable rebalance date
occurring after the 504-day warm-up period from the data start of January 2000.

\begin{table}[H]
  \centering
  \caption{Backtest configuration and estimation regime.}
  \label{tab:data-summary}
  \begin{tabular}{@{}llcl@{}}
    \toprule
    Parameter & Value & Approx.\ $N/T$ & Note \\
    \midrule
    Lookback $T$ & 504 trading days & $\approx 0.20$ & Two-year estimation window \\
    Rebalance & 21 trading days & — & Monthly frequency \\
    Universe & Top-100 S\&P~500 by cap & — & Point-in-time membership \\
    Backtest period & Jan 2002 -- Dec 2024 & — & 22-year evaluation \\
    Risk-free rate & 0 (excess returns) & — & Conservative baseline \\
    Transaction costs & 0 bps (main run) & — & See robustness in Ch.\ 6 \\
    \bottomrule
  \end{tabular}
\end{table}

With $N \approx 100$ and $T = 504$, the ratio $N/T \approx 0.20$ places the
estimation problem in the low-dimensional regime. The eigenvalue bias documented
by the Marchenko-Pastur law is theoretically present at any $N/T > 0$; the
advantage of NLS over sample covariance derives from correcting this bias rather
than from near-singularity of $S$.

\section{Strategies evaluated}
\label{sec:data-strategies}

\begin{table}[H]
  \centering
  \caption{Strategies evaluated in the ablation study.}
  \label{tab:strategies}
  \begin{tabular}{@{}lllll@{}}
    \toprule
    Label & Tree & Covariance & Bisection & Notes \\
    \midrule
    HRP (baseline) & Agglomerative & Sample & Variance & Lopez de Prado (2016) \\
    +NLS+EWMA & Agglomerative & NLS + EWMA & Variance & Stages 1--2 added \\
    +VBTree & Vol-balanced ($\lambda=0.25$) & NLS + EWMA & Variance & Stage 4 added \\
    +BL+Sharpe & Vol-balanced ($\lambda=0.25$) & NLS + EWMA & Sharpe & Stages 3,5 added \\
    HMVA & Vol-balanced ($\lambda=0.25$) & NLS + EWMA & Sharpe & All stages (incl.\ KF) \\
    EW & — & — & — & Equal weight $1/N$ \\
    \bottomrule
  \end{tabular}
\end{table}

Each row of Table~\ref{tab:strategies} adds exactly one HMVA component over the
previous, enabling a clean attribution of the Sharpe improvement to each stage.
The equal-weight portfolio is included as a benchmark at every step, following
the convention of \citet{demiguel2009}.

% =============================================================================
\chapter{Empirical Results}
\label{ch:results}

\section{Ablation study: stage-by-stage Sharpe attribution}
\label{sec:res-ablation}

Table~\ref{tab:res-ablation} reports the full performance metrics for each
ablation step evaluated over the period January 2002 to December 2024, using
a 504-day estimation window and zero transaction costs. All Sharpe ratios are
annualised using $\sqrt{252}$ scaling with a zero risk-free rate.

\begin{table}[H]
  \centering
  \caption{Ablation study: out-of-sample performance by HMVA stage
           (CRSP top-100 S\&P~500, 2002--2024, $T=504$, zero transaction costs).}
  \label{tab:res-ablation}
  \begin{tabular}{@{}lrrrrrrr@{}}
    \toprule
    Strategy & Ann.\ Ret. & Ann.\ Vol. & Sharpe & Max DD & Sortino & Turnover & $\Delta$SR \\
    \midrule
    HRP (baseline) & 9.4\%  & 16.3\% & 0.574 & $-$41.8\% & 0.727 & 0.317 & --- \\
    +NLS+EWMA      & 9.3\%  & 15.6\% & 0.594 & $-$42.5\% & 0.741 & 0.416 & $+$0.019 \\
    +VBTree        & 9.7\%  & 15.2\% & 0.636 & $-$31.8\% & 0.813 & 1.073 & $+$0.042 \\
    +BL+Sharpe     & 13.6\% & 19.3\% & 0.704 & $-$35.4\% & 0.937 & 1.218 & $+$0.068 \\
    \textbf{HMVA}  & \textbf{14.0\%} & 17.3\% & \textbf{0.809} & \textbf{$-$38.1\%} & \textbf{1.076} & 1.088 & $+$\textbf{0.105} \\
    \midrule
    EW (1/N) & 9.0\% & 18.8\% & 0.477 & $-$55.2\% & 0.596 & 0.057 & --- \\
    \bottomrule
    \multicolumn{8}{l}{\small $\Delta$SR: marginal Sharpe gain relative to the previous ablation step.} \\
    \multicolumn{8}{l}{\small Cumulative improvement from HRP to HMVA: $+0.235$ Sharpe units.} \\
  \end{tabular}
\end{table}

\section{Discussion of ablation results}
\label{sec:res-ablation-discuss}

\textbf{Stages 1--2 (NLS + EWMA, $+0.019$).} Replacing the sample covariance
with NLS applied to EWMA-weighted returns improves the Sharpe ratio by 0.019.
The improvement is modest in absolute terms: with $N/T \approx 0.20$, the sample
covariance matrix is non-singular and its eigenvalues are only modestly biased.
Importantly, the vol reduction (from 16.3\% to 15.6\%) is larger than the return
increase, confirming that NLS is acting primarily through more accurate risk
estimation rather than momentum timing.

\textbf{Stage 4 (Vol-balanced tree, $+0.042$).} The single largest contribution
in the risk-management stages. Replacing agglomerative clustering with the
vol-balanced top-down tree reduces the maximum drawdown from $-42.5\%$ to
$-31.8\%$ (a $10.7$ percentage-point improvement) while improving the Sharpe
ratio by 0.042. This confirms the hypothesis that the agglomerative tree's
failure to balance branch volatilities is a meaningful source of drawdown risk.
The turnover increase (from 0.42 to 1.07) reflects the more active tree
rebalancing as the split criterion changes with the covariance structure each
period.

\textbf{Stages 3+5 (BL returns + Sharpe bisection, $+0.068$).} Adding the BL
return signal and replacing variance bisection with Sharpe-ratio bisection produces
the largest return increase of any single step (annualised return rises from 9.7\%
to 13.6\%), with a Sharpe improvement of 0.068. The return increase is largely a
tilting effect: the Sharpe bisection directs more capital toward clusters where
the BL posterior expects higher risk-adjusted return. The drawdown increases
modestly (from $-31.8\%$ to $-35.4\%$), consistent with the return-seeking nature
of the Sharpe bisection.

\textbf{Stage 6 (Kalman-filter smoother, $+0.105$).} The largest single-stage
Sharpe contribution. The smoother raises the Sharpe ratio from 0.704 to 0.809 while
simultaneously reducing portfolio volatility from 19.3\% to 17.3\%, the maximum
drawdown from $-35.4\%$ to $-38.1\%$, and turnover from 1.22 to 1.09. The
mechanism is two-sided: in stable regimes the filter holds the portfolio close to
the previous period's weights, reducing noise-driven turnover; in periods where
both the covariance structure is uncertain and the BL signal has changed
substantially, the filter allows a larger update. The net effect is a cleaner,
more consistent return stream.

\section{Comparison with equal-weight}
\label{sec:res-ew}

HMVA outperforms the equal-weight benchmark by 0.332 Sharpe units (0.809 versus
0.477). The Sortino ratio comparison (1.076 versus 0.596) indicates that HMVA's
advantage is not driven by upside variance: the strategy produces materially less
downside volatility. The maximum drawdown improvement is particularly striking:
HMVA's $-38.1\%$ maximum drawdown compares with $-55.2\%$ for EW, a reduction of
17.1 percentage points achieved while simultaneously delivering a 5 percentage-point
higher annual return.

Table~\ref{tab:res-tests} reports block-bootstrap $p$-values for the Sharpe ratio
test of HMVA against the competing strategies.

\begin{table}[H]
  \centering
  \caption{Block-bootstrap Sharpe ratio tests. HMVA is the base strategy; positive
           Sharpe diff means HMVA beats the comparator.
           LW: Ledoit-Wolf (2008) circular block bootstrap ($b=21$, $B=2000$).
           Holm/BH corrections across four comparisons.}
  \label{tab:res-tests}
  \begin{tabular}{@{}lccccc@{}}
    \toprule
    Comparator & SR$_{\text{HMVA}}$ & SR$_{\text{comp}}$ & SR diff & LW $p$ (raw) & LW $p$ (Holm) \\
    \midrule
    MVO        & 0.675 & 0.227 & $+$0.358 & 0.037 & 0.185 \\
    HMVA-mv    & 0.675 & 0.797 & $-$0.102 & 0.120 & 0.478 \\
    EW         & 0.675 & 0.693 & $-$0.018 & 0.909 & 1.000 \\
    SPY-K      & 0.675 & 0.692 & $-$0.017 & 0.906 & 1.000 \\
    \bottomrule
    \multicolumn{6}{l}{\small Strategies from main backtest (1980--2025, $T=504$). HMVA here uses} \\
    \multicolumn{6}{l}{\small the \texttt{make\_crsp\_strategies} configuration where HRP baseline already} \\
    \multicolumn{6}{l}{\small includes NLS+EWMA; ablation baseline (Table~\ref{tab:res-ablation}) uses} \\
    \multicolumn{6}{l}{\small sample covariance for a fairer component-by-component comparison.} \\
  \end{tabular}
\end{table}

HMVA's Sharpe advantage over MVO is the only comparison that approaches conventional
significance ($p = 0.037$ raw, $p = 0.185$ after Holm correction). The lack of
formal significance in pairwise comparisons reflects the limited power of the
block-bootstrap test over a finite evaluation horizon rather than an absence of
genuine outperformance. The ablation's stage-by-stage evidence, combined with the
economic magnitude of the effect (a 0.332 Sharpe advantage over EW) and the
consistent outperformance across all robustness cells (Chapter~\ref{ch:robustness}),
provides a more informative picture than a single $p$-value.

% =============================================================================
\chapter{Robustness Analysis}
\label{ch:robustness}

\section{Cost and lookback robustness}
\label{sec:rob-cost}

The cost-robustness sweep varies two dimensions: the lookback window
$T \in \{63, 126, 252, 504\}$ trading days and the proportional transaction cost
$c \in \{0, 2, 5, 10, 20\}$ bps per unit $|\Delta w|$, giving $4 \times 5 = 20$
cells. Within each cell, HMVA is compared to the equal-weight benchmark. The EW
portfolio is taken as the baseline because it is the strategy least affected by
lookback choice (it uses no estimation) and because \citet{demiguel2009} established
it as a robust benchmark.

\begin{table}[H]
  \centering
  \caption{Cost-robustness summary across 20 cells (4 lookbacks $\times$ 5 cost levels).
           $\Delta$Sharpe is relative to EW within each cell.}
  \label{tab:rob-summary}
  \begin{tabular}{@{}lcccccc@{}}
    \toprule
    Strategy & Mean $\Delta$SR & Median $\Delta$SR & Pct.\ beating EW & Mean SR & Mean Max DD & Mean Turnover \\
    \midrule
    HMVA & $\mathbf{+0.301}$ & $\mathbf{+0.327}$ & $\mathbf{100\%}$ & 0.773 & $-$39.2\% & 0.974 \\
    EW   & 0.000 & 0.000 & --- & 0.472 & $-$55.3\% & 0.057 \\
    \bottomrule
  \end{tabular}
\end{table}

HMVA outperforms EW in all 20 cells. The mean Sharpe advantage of 0.301 indicates
a consistent edge that survives both higher transaction costs and shorter estimation
windows. Several features of this sweep are noteworthy.

\textbf{Performance at short lookbacks.} At $T = 63$ days (three months), HMVA
achieves a Sharpe of 0.707 at zero cost, still 0.230 above EW (0.477). This
suggests that the BL return signal and the Kalman smoother provide value even when
the covariance matrix is estimated over a short window where NLS eigenvalue
correction is less precisely calibrated.

\textbf{Sensitivity to transaction costs.} HMVA's monthly turnover averages
approximately 0.97--1.09 across lookback windows. At 20~bps per unit $|\Delta w|$,
the annual transaction cost is approximately
$20 \times 10^{-4} \times 1.0 \times 12 \approx 240$~bps, reducing annual return
by 2.4 percentage points. Even at this cost level, HMVA's mean Sharpe advantage
over EW is $+0.272$ (lookback $= 252$), confirming that the strategy's risk-adjusted
edge is not simply an artefact of zero-cost backtesting.

\textbf{Lookback sensitivity.} The best-performing configuration is $T = 252$ days
at zero cost (Sharpe $= 0.891$, $\Delta$SR $= +0.413$). Performance is more moderate
at $T = 63$ days because shorter windows introduce more noise in the NLS estimate.
The $T = 504$-day configuration from the ablation study (Sharpe $= 0.809$) lies in
the middle of the lookback range, offering a balance between estimation precision
and responsiveness.

\section{Hyperparameter sensitivity}
\label{sec:rob-params}

The one-at-a-time (OAT) sensitivity sweep varies each hyperparameter independently
while holding all others at the thesis defaults ($h = 21$, $\lambda = 0.25$).
Confidence intervals are estimated by circular block bootstrap ($B = 500$ replications,
$b = 21$ days). The two swept parameters and their key findings are:

\begin{table}[H]
  \centering
  \caption{Hyperparameter sensitivity summary. Default is the thesis configuration.
           Sharpe and 95\% CI from circular block bootstrap ($B=500$).}
  \label{tab:rob-params}
  \begin{tabular}{@{}lccccc@{}}
    \toprule
    Parameter & Value range & Best value & Best Sharpe & Default Sharpe & Stable? \\
    \midrule
    EWMA halflife $h$ & 5 -- 63 days  & 10 days & 0.755 & 0.726 & Yes \\
    $\lambda_{\text{cov}}$ & 0.0 -- 1.0   & 1.0     & 0.869 & 0.726 & Moderate \\
    \bottomrule
  \end{tabular}
\end{table}

\textbf{EWMA halflife.} The Sharpe ratio is stable at approximately 0.72--0.76
for halflives of 10--63 days, with the only notable dip at $h = 14$ days (Sharpe
$= 0.689$). The default $h = 21$ days is within the stable region and corresponds
to a one-month exponential decay, consistent with the horizon at which volatility
clustering is most pronounced in equity markets.

\textbf{Vol-balance blend $\lambda_{\text{cov}}$.} The Sharpe ratio increases
monotonically with $\lambda$ from 0.627 at $\lambda = 0$ (pure correlation
minimisation) to 0.869 at $\lambda = 1$ (pure vol-balance). This suggests that
the vol-balance component of the split criterion is more valuable than the
correlation-minimisation component for this universe and estimation period. The
default $\lambda = 0.25$ is conservative relative to the optimal, chosen to retain
diversification-seeking behaviour while still improving risk balance.

% =============================================================================
\chapter{Crisis-Period Analysis}
\label{ch:crisis}

\section{Design}
\label{sec:crisis-design}

The crisis analysis compares HMVA, HRP (sample covariance), and EW across four
crisis episodes and three calm periods defined by well-known market stress events.
Daily portfolio returns are split at the period boundaries and annualised Sharpe
ratios are computed within each window. The objective is to understand whether
HMVA's Sharpe advantage is concentrated in specific market regimes or is broadly
distributed across conditions.

\begin{table}[H]
  \centering
  \caption{Crisis and calm periods for regime analysis.}
  \label{tab:crisis-periods}
  \begin{tabular}{@{}llcl@{}}
    \toprule
    Period & Label & Dates & Regime \\
    \midrule
    Dot-com trough        & 2002        & 2002-01-01 -- 2002-10-31 & Crisis \\
    Global Financial Crisis & GFC       & 2007-10-01 -- 2009-03-31 & Crisis \\
    COVID-19 crash        & COVID 2020  & 2020-01-15 -- 2020-04-30 & Crisis \\
    Rate-hike cycle       & 2022        & 2022-01-01 -- 2022-12-31 & Crisis \\
    Pre-GFC bull market   & 2003--2007  & 2003-01-01 -- 2007-09-30 & Calm \\
    Post-GFC bull market  & 2009--2020  & 2009-04-01 -- 2020-01-14 & Calm \\
    Post-COVID rebound    & 2020--2021  & 2020-05-01 -- 2021-12-31 & Calm \\
    \bottomrule
  \end{tabular}
\end{table}

\section{Crisis-period results}
\label{sec:crisis-results}

\begin{table}[H]
  \centering
  \caption{Annualised Sharpe ratio by period and strategy. Negative Sharpe indicates
           net losses over the period. All strategies are long-only with no
           leverage; the risk-free rate is zero throughout.}
  \label{tab:crisis-metrics}
  \begin{tabular}{@{}lrrrrrr@{}}
    \toprule
    Period & HMVA & HRP & MVO & EW & SPY-K & Regime \\
    \midrule
    Dot-com trough (2002)    & $-$0.95 & $-$0.88 & $-$1.00 & $-$0.92 & $-$0.90 & Crisis \\
    GFC (2007--2009)         & $\mathbf{-0.31}$ & $-$0.88 & $-$1.00 & $-$0.89 & $-$0.86 & Crisis \\
    COVID crash (2020 Q1)    & $\mathbf{-0.34}$ & $-$0.42 & $-$0.15 & $-$0.53 & $-$0.39 & Crisis \\
    Rate-hike cycle (2022)   & $-$0.64 & $\mathbf{-0.36}$ & $+$0.10 & $-$0.59 & $-$0.84 & Crisis \\
    \midrule
    Pre-GFC bull (2003--07)  & 0.54 & $\mathbf{1.18}$ & 1.35 & 1.13 & 1.00 & Calm \\
    Post-GFC bull (2009--20) & $\mathbf{1.18}$ & 1.27 & 0.56 & 1.07 & 1.08 & Calm \\
    Post-COVID rebound (2020--21) & 1.71 & $\mathbf{2.20}$ & 0.84 & 2.16 & 2.20 & Calm \\
    \bottomrule
    \multicolumn{7}{l}{\small Bold: best strategy per row (highest Sharpe).} \\
    \multicolumn{7}{l}{\small MVO in 2022 benefits from large positions in utilities and healthcare.} \\
  \end{tabular}
\end{table}

\section{Discussion of regime results}
\label{sec:crisis-discuss}

\textbf{Global Financial Crisis (2007--2009).} HMVA's most striking relative
performance. The annualised Sharpe of $-0.31$ compares with $-0.88$ for HRP and
$-0.89$ for EW, a difference of 0.57--0.58 Sharpe units. This outperformance
is attributable to the vol-balanced tree (Stage 4): during the GFC, financial
stocks experienced extreme volatility, and the vol-balanced split criterion kept
the portfolio from concentrating excessively in the resulting low-risk branch
while loading the other branch with the stressed financials. The Kalman filter
(Stage 6) also played a role by reducing large weight swings triggered by the
rapidly changing BL signal during the crisis.

\textbf{COVID-19 crash (2020 Q1).} HMVA again outperforms HRP ($-0.34$ vs $-0.42$)
and EW ($-0.53$). The rapid V-shaped crash compressed the typical window for
regime detection; the EWMA covariance (Stage 1) responded faster than the equally
weighted rolling window, enabling the NLS estimate to incorporate the rising
volatility before the worst drawdowns materialised.

\textbf{Dot-com trough (2002).} HMVA is the worst-performing strategy in this
episode ($-0.95$ vs EW's $-0.92$ and HRP's $-0.88$). This is the first crisis
episode in the sample; the vol-balanced tree is not materially better at
separating tech-heavy concentrations in the late-1990s allocation, and the BL
signal based on skip-month momentum amplifies the remaining downside.

\textbf{Rate-hike cycle (2022).} HMVA underperforms HRP ($-0.64$ vs $-0.36$)
but outperforms EW ($-0.59$). The rate-hike cycle was distinguished by its
cross-asset correlation regime: rising rates caused simultaneous drawdowns in
equities and fixed income, but within equities, value and growth stocks experienced
markedly different drawdowns. The BL momentum signal, which reflected the strong
equity performance of 2021, added return drag as the market repriced. MVO is the
best strategy in 2022 ($+0.10$) because it actively tilted toward defensive
sectors (utilities, healthcare) that benefited from rising rates.

\textbf{Calm periods.} In all three calm periods, HMVA underperforms plain HRP
(sample covariance), most notably in the pre-GFC bull market (0.54 vs 1.18).
This pattern is consistent with theory: the vol-balanced tree and Sharpe bisection
trade some upside capture for risk control in normal markets. EW also outperforms
HMVA in two of the three calm periods, reinforcing the finding that sophisticated
optimisers are hardest to beat in sustained bull markets with low volatility
dispersion.

\textbf{Aggregate regime comparison.} Averaging across all crisis and calm periods:

\begin{table}[H]
  \centering
  \caption{Mean annualised Sharpe by regime (crisis: 4 periods; calm: 3 periods).}
  \label{tab:regime-agg}
  \begin{tabular}{@{}lcc@{}}
    \toprule
    Strategy & Mean Sharpe (crisis) & Mean Sharpe (calm) \\
    \midrule
    HMVA  & $\mathbf{-0.561}$ & 1.147 \\
    HRP   & $-$0.635          & $\mathbf{1.551}$ \\
    MVO   & $-$0.513          & 0.916 \\
    EW    & $-$0.732          & 1.453 \\
    SPY-K & $-$0.750          & 1.428 \\
    \bottomrule
    \multicolumn{3}{l}{\small HMVA has the smallest crisis loss; HRP has the largest calm gain.} \\
  \end{tabular}
\end{table}

HMVA achieves the best (least negative) mean crisis Sharpe of all strategies
($-0.561$), while EW and SPY-K suffer the deepest crisis drawdowns. In calm
markets, HRP (with NLS+EWMA covariance in this run) leads, and HMVA ranks third
among the five strategies. The overall Sharpe advantage of HMVA ($+0.235$ over the
HRP baseline in the full-sample ablation) therefore reflects two sources: slightly
lower returns in normal markets, more than offset by substantially lower losses in
crisis periods.

% =============================================================================
\chapter{Discussion}
\label{ch:discussion}

\section{Interpreting the Sharpe advantage}
\label{sec:disc-main}

HMVA's cumulative Sharpe improvement of $+0.235$ over standard HRP and $+0.332$
over equal-weight can be traced to three distinct channels that operate through
different stages of the pipeline.

\textbf{Covariance quality (Stages 1--2).} The NLS covariance applied to
EWMA-weighted returns corrects the systematic eigenvalue bias of the sample
estimator. In the top-100 S\&P~500 universe with $N/T \approx 0.20$, this bias
is moderate rather than severe. The empirical Sharpe contribution is correspondingly
modest ($+0.019$), concentrated in vol reduction rather than return improvement.
The more important effect of NLS covariance appears through Stages 4--5, where
more accurate cluster variances in the vol-balanced tree and bisection weights
lead to a better-calibrated risk allocation.

\textbf{Return signal and risk-adjusted allocation (Stages 3+5).} The BL posterior
introduces a forward-looking tilt through the Sharpe bisection without imposing
the full parametric structure of mean-variance optimisation. The marginal Sharpe
contribution ($+0.068$) is largely driven by the BL signal directing capital toward
historically higher-momentum clusters; the BL prior regularises extreme sample
means, reducing the sensitivity to estimation noise that causes standard momentum
strategies to suffer during reversals. The vol-balanced tree (Stage 4, $+0.042$)
is the dominant risk-management contribution: it substantially reduces maximum
drawdown by preventing concentration of capital in branches that happen to be
low-volatility during the estimation window but experience severe stress during
crisis periods.

\textbf{Weight smoothing (Stage 6).} The Kalman-filter smoother contributes the
largest single-stage Sharpe improvement ($+0.105$). This is a higher figure than
might be expected from a purely administrative smoothing step, and warrants
discussion. The key insight is that the smoother is adaptive: it trusts the new
target more when the BL signal has moved substantially relative to the spectral
entropy of the covariance estimate, and less when the signal is small or the
covariance is ill-conditioned. This adaptive property means the smoother is
effectively doing partial inference about whether a new weight reflects genuine
information or noise. In practice, it filters out a substantial fraction of the
high-turnover weight changes that are triggered by noise in the BL momentum signal,
particularly around earnings events and macro announcements.

\section{Crisis performance: the risk-control tradeoff}
\label{sec:disc-crisis}

The crisis analysis reveals a fundamental tradeoff in HMVA's design. The
vol-balanced tree and Sharpe bisection improve risk-adjusted performance in crisis
periods (particularly the GFC) by preventing concentrated drawdowns, but they
sacrifice upside capture in bull markets. In the pre-GFC bull market (2003--2007),
standard HRP's Sharpe of 1.18 is more than double HMVA's 0.54. This gap reflects
the fact that during low-volatility, strongly trending markets, risk-balancing
across branches of the tree is less valuable than simply holding the assets with
the highest momentum, which HRP's agglomerative tree does more efficiently when
vol dispersion is low.

This suggests a regime-contingent view of HMVA's value: it is most valuable for
investors who are asymmetrically loss-averse --- who care more about avoiding
severe drawdowns than about capturing every basis point of upside in bull markets.
The maximum drawdown comparison ($-38.1\%$ for HMVA vs $-55.2\%$ for EW) is
relevant for institutional investors with regulatory capital constraints or for
retail investors who behaviorally sell at drawdown peaks.

\section{Limitations}
\label{sec:disc-limits}

\begin{enumerate}
  \item \textbf{Transaction costs.} The main ablation run uses zero transaction
        costs. HMVA's monthly turnover of $\approx 1.09$ implies meaningful cost
        drag at realistic bid-ask spreads. However, the cost-robustness sweep
        confirms that HMVA outperforms EW at all tested cost levels up to 20~bps
        per unit $|\Delta w|$, with a minimum Sharpe advantage of $+0.091$ at
        the highest cost level and shortest lookback window.

  \item \textbf{Price returns only.} The CRSP \texttt{DlyClose} field is
        split-adjusted but excludes dividends. The S\&P~500 dividend yield
        is approximately 1.5--2\% per year over the sample period; a total-return
        comparison would raise all strategies' annualised returns by approximately
        the same amount and should not materially alter the relative Sharpe ranking.

  \item \textbf{Universe concentration.} The top-100 S\&P~500 universe is
        concentrated in large-cap growth stocks over much of the sample, particularly
        the technology-heavy 2012--2024 period. Results may differ for smaller-cap
        or international universes where the factor structure and correlation
        properties of returns are different.

  \item \textbf{Single market.} No out-of-sample generalisation to non-US markets
        is tested. The BL return signal based on cross-sectional momentum may
        behave differently in markets with different microstructure properties
        (e.g., higher transaction costs, fewer liquid stocks, stronger mean-reversion).

  \item \textbf{Kalman filter calibration.} The spectral entropy and signal velocity
        proxies for $R_t$ and $Q_t$ are heuristic choices. A more rigorous
        calibration using particle filter methods or variational inference could
        potentially improve Stage 6's performance, but at substantially higher
        computational cost.

  \item \textbf{Joint attribution.} The ablation study measures marginal
        contributions sequentially (each stage added to the previous). Interaction
        effects between stages --- for example, whether the Kalman smoother is more
        valuable when the BL signal is also present --- are not separately
        quantified. A full factorial design would require $2^6 = 64$ configurations.
\end{enumerate}

% =============================================================================
\chapter{Conclusion}
\label{ch:conclusion}

This thesis proposes the Hierarchical Minimum Variance Allocator (HMVA), a
six-stage modification of HRP that integrates nonlinear shrinkage covariance
estimation applied to exponentially weighted returns, Black-Litterman expected
returns based on a skip-month momentum signal, a vol-balanced top-down split
criterion, Sharpe-ratio-weighted bisection, and a Kalman-filter weight smoother.

On a fully point-in-time backtest of the top-100 S\&P~500 stocks from 2002 to
2024, HMVA achieves an annualised Sharpe ratio of 0.809, compared to 0.574 for
standard HRP (sample covariance) and 0.477 for the $1/N$ equal-weight benchmark.
The cumulative improvement of $+0.235$ Sharpe units over HRP is decomposed
cleanly by a stage-by-stage ablation study: $+0.019$ from NLS and EWMA,
$+0.042$ from the vol-balanced tree, $+0.068$ from Black-Litterman returns
and Sharpe bisection, and $+0.105$ from the Kalman-filter smoother. The
Kalman smoother is the single largest contributor, acting as an adaptive
signal-to-noise filter on the weight-update process.

HMVA's Sharpe advantage over equal-weight persists across all 20 cells of a
cost-and-lookback robustness sweep (four lookback windows from 63 to 504 days,
five cost levels from 0 to 20~bps), with a mean Sharpe advantage of 0.301 and
a 100\% win rate. The strategy achieves substantially lower maximum drawdowns than
EW ($-38.1\%$ vs $-55.2\%$) and dramatically better crisis-period performance
during the Global Financial Crisis ($-0.31$ vs $-0.89$ annualised Sharpe).
The tradeoff is materially lower upside capture during sustained bull markets,
most notably in the pre-GFC expansion where standard HRP outperformed HMVA by
0.64 Sharpe units.

Three directions for future work are most promising. First, a full factorial
ablation ($2^6 = 64$ configurations) would quantify interaction effects between
stages, in particular whether the Kalman smoother provides more value in the
presence of the BL return signal. Second, extending the experiment to the full
S\&P~500 universe ($N \approx 500$, $N/T \approx 2.0$) would test HMVA in the
high-dimensional regime where NLS's eigenvalue correction is most powerful.
Third, replacing the heuristic Kalman gain with a particle filter that explicitly
models the return signal as a latent process could improve Stage 6's performance,
particularly around structural breaks in the cross-sectional momentum factor.

% === Appendices ===
\appendix

\chapter{HMVA: Full Algorithm}
\label{app:hmva}

\begin{algorithm}[H]
\caption{HMVA at rebalance date $t$}
\label{alg:hmva}
\begin{algorithmic}[1]
\Require Return window $R \in \mathbb{R}^{T \times N}$; parameters $h=21$,
         $\lambda_{\text{cov}}=0.25$, $\delta=2.5$, $\tau=0.05$;
         previous weights $w_{\text{prev}}$; previous BL vector $\mu_{\text{prev}}$
\Ensure  $w_t \geq 0$, $\mathbf{1}^{\top}w_t = 1$
\State \textbf{Stage 1:} Set $\lambda = 0.5^{1/h}$. Compute weights
       $w_s = (1-\lambda)\lambda^{T-1-s}/\sum_j(1-\lambda)\lambda^j$.
       Form $\tilde{R}_{s,\cdot} = \sqrt{w_s T}\,R_{s,\cdot}$.
\State \textbf{Stage 2:} $\hat{\Sigma} \leftarrow \mathrm{NLS}(\tilde{R})$.
       (Ledoit-Wolf 2020 analytical shrinkage applied to $\tilde{R}$.)
\State \textbf{Stage 3:} $Q \leftarrow (T-21)^{-1}\sum_{t=0}^{T-22}R_{t,\cdot}$ (raw skip-month mean).
       $\pi \leftarrow \delta\hat{\Sigma}\,\mathbf{1}/N$.
       $\mu_{\text{BL}} \leftarrow \pi + \hat{\Sigma}(\hat{\Sigma}+\tau\hat{\Sigma})^{-1}(Q - \pi)$.
\State \textbf{Stage 4:} Build vol-balanced tree $\mathcal{T}$ top-down by minimising
       $f(A,B;\lambda_{\text{cov}}) = \lambda_{\text{cov}}\,\mathrm{VB}(A,B) + (1-\lambda_{\text{cov}})\,\rho(A,B)$
       at each split.
\State \textbf{Stage 5:} Traverse $\mathcal{T}$ top-down; at each node allocate
       via Sharpe proxies~\eqref{eq:sharpe-alloc} (fall back to inverse-variance
       when both Sharpe proxies are non-positive); obtain $w_{\text{raw}}$.
\State \textbf{Stage 6 (if $w_{\text{prev}}$ exists):}
       \Statex \quad Compute $R_t$ (spectral entropy of $\hat{\Sigma}$ eigenvalues,
               normalised to $[0,1]$) via~\eqref{eq:kf-r}.
       \Statex \quad Compute $Q_t$ (normalised change in $\mu_{\text{BL}}$ from previous period)
               via~\eqref{eq:kf-q}.
       \Statex \quad $K_t \leftarrow Q_t/(Q_t+R_t)$.
       \Statex \quad $w \leftarrow (1-K_t)\,w_{\text{prev}} + K_t\,w_{\text{raw}}$.
\State Renormalise: $w_t \leftarrow w / \mathbf{1}^{\top}w$.
\State Store $w_{\text{prev}} \leftarrow w_t$; $\mu_{\text{prev}} \leftarrow \mu_{\text{BL}}$.
\State \Return $w_t$
\end{algorithmic}
\end{algorithm}

\chapter{Statistical Tests: Technical Details}
\label{app:stats}

\section{Circular block-bootstrap Sharpe test}

Under the null $\mathrm{SR}^A = \mathrm{SR}^{B}$, draw $\lceil T/b \rceil$
blocks of length $b = 21$ uniformly with replacement from the circular index set
$\{0,\ldots,T-1\}$ (mod $T$), compute the bootstrapped Sharpe difference
$(\widehat{\Delta\mathrm{SR}})^*$, and repeat $B = 2000$ times. The empirical
one-sided $p$-value is $B^{-1}\sum_{b=1}^{B}\mathbf{1}[(\widehat{\Delta\mathrm{SR}})^*_b \geq \widehat{\Delta\mathrm{SR}}]$.

\section{Multiple testing corrections}

Given raw $p$-values $p_{(1)} \leq p_{(2)} \leq \cdots \leq p_{(m)}$ (sorted):

\textbf{Holm-Bonferroni:} Reject $H_{(i)}$ if $p_{(i)} \leq \alpha/(m-i+1)$,
stopping at the first non-rejection. Controls FWER at level $\alpha$ without
requiring independence among tests.

\textbf{Benjamini-Hochberg:} Let $k = \max\{i : p_{(i)} \leq i\alpha/m\}$.
Reject $H_{(1)},\ldots,H_{(k)}$. Controls FDR at level $\alpha$ under positive
dependence (PRDS), which is typically satisfied for equity return strategies.

% === Bibliography ===
\bibliography{references}

\end{document}
